% !TeX program = lualatex
% !TeX root = predavanje00.tex
% ============================================================
% Zajednički preamble za sva predavanja RWA kolegija
% Uključiti s: % ============================================================
% Zajednički preamble za sva predavanja RWA kolegija
% Uključiti s: % ============================================================
% Zajednički preamble za sva predavanja RWA kolegija
% Uključiti s: \input{../common/preamble.tex}
% ============================================================

\documentclass[aspectratio=149,11pt]{beamer}
\usetheme{CambridgeUS}

\usepackage{fontspec}              % radi samo sa XeLaTeX/LuaLaTeX
\setsansfont{Latin Modern Sans}    % sans-serif font
\setmonofont{Latin Modern Mono}    % monospace font za kod
\usepackage[croatian]{babel}
\usepackage{csquotes}              % \enquote{} → hrvatski navodnici „..."
\usepackage{lmodern}
\usepackage{booktabs}
\usepackage{listings}
\usepackage{xcolor}
\usepackage{hyperref}

% ---------- Autor i logo ----------
\author{mag.ing. Josip Torić (SIT UNIZD)}
\titlegraphic{\includegraphics[height=2.5cm]{../common/SIT_logo_tekst1.pdf}}

% ---------- Listings ----------
\definecolor{codebg}{RGB}{248,248,248}
\definecolor{codeframeclr}{RGB}{220,220,220}
\lstdefinelanguage{bash}{
  morekeywords={sudo,apt,brew,pip,pipx,uv,winget,python,git,flask,source,deactivate,code,gunicorn,docker,compose,heroku,render,railway,node,npm,pnpm,yarn,uvicorn,alembic,pytest,ruff,black,precommit},
  sensitive=false,
  morecomment=[l]{\#}
}
\lstdefinelanguage{python}{
  morekeywords={from,import,def,class,return,if,elif,else,for,in,while,try,except,raise,with,as,not,and,or,is,lambda,yield,True,False,None,await,async},
  sensitive=true,
  morecomment=[l]{\#},
  morestring=[b]',
  morestring=[b]"
}
\lstdefinelanguage{sql}{
  morekeywords={SELECT,FROM,WHERE,INSERT,INTO,VALUES,UPDATE,SET,DELETE,CREATE,TABLE,PRIMARY,KEY,FOREIGN,REFERENCES,UNIQUE,NOT,NULL,INDEX,ON,ALTER,DROP,CONSTRAINT},
  sensitive=false,
  morecomment=[l]{--}
}
\lstdefinelanguage{json}{
  morecomment=[l]{//},
  morestring=[b]",
  showstringspaces=false
}
\lstdefinelanguage{typescript}{
  morekeywords={import,export,from,const,let,var,function,return,if,else,for,of,in,while,switch,case,break,async,await,try,catch,finally,throw,new,class,this,extends,implements,interface,type,enum,readonly,public,private,protected,static,get,set,true,false,null,undefined,void,string,number,boolean,any,unknown,never,typeof,keyof,as},
  sensitive=true,
  morecomment=[l]{//},
  morecomment=[s]{/*}{*/},
  morestring=[b]',
  morestring=[b]",
  morestring=[b]`
}
\lstalias{javascript}{typescript}
\lstset{
  backgroundcolor=\color{codebg},
  basicstyle=\ttfamily\footnotesize,
  frame=single,
  rulecolor=\color{codeframeclr},
  frameround=tttt,
  breaklines=true,
  tabsize=2,
  showstringspaces=false,
  upquote=true,
  columns=fullflexible
}

% ---------- Helper macros ----------
\newcommand{\codefile}[1]{\texttt{#1}}
\newcommand{\kbd}[1]{\texttt{#1}}

% ---------- Code frame environment ----------
% Automatski dodaje fragile i shrink=5, koristi se umjesto \begin{frame}
% Primjer: \begin{codeframe}{Naslov} ... \end{codeframe}
% Primjer s opcijama: \begin{codeframe}[shrink=15]{Naslov} ... \end{codeframe}
\newenvironment{codeframe}[2][]
  {\begin{frame}[fragile,shrink=5,#1]{#2}}
  {\end{frame}}

% ============================================================

\documentclass[aspectratio=149,11pt]{beamer}
\usetheme{CambridgeUS}

\usepackage{fontspec}              % radi samo sa XeLaTeX/LuaLaTeX
\setsansfont{Latin Modern Sans}    % sans-serif font
\setmonofont{Latin Modern Mono}    % monospace font za kod
\usepackage[croatian]{babel}
\usepackage{csquotes}              % \enquote{} → hrvatski navodnici „..."
\usepackage{lmodern}
\usepackage{booktabs}
\usepackage{listings}
\usepackage{xcolor}
\usepackage{hyperref}

% ---------- Autor i logo ----------
\author{mag.ing. Josip Torić (SIT UNIZD)}
\titlegraphic{\includegraphics[height=2.5cm]{../common/SIT_logo_tekst1.pdf}}

% ---------- Listings ----------
\definecolor{codebg}{RGB}{248,248,248}
\definecolor{codeframeclr}{RGB}{220,220,220}
\lstdefinelanguage{bash}{
  morekeywords={sudo,apt,brew,pip,pipx,uv,winget,python,git,flask,source,deactivate,code,gunicorn,docker,compose,heroku,render,railway,node,npm,pnpm,yarn,uvicorn,alembic,pytest,ruff,black,precommit},
  sensitive=false,
  morecomment=[l]{\#}
}
\lstdefinelanguage{python}{
  morekeywords={from,import,def,class,return,if,elif,else,for,in,while,try,except,raise,with,as,not,and,or,is,lambda,yield,True,False,None,await,async},
  sensitive=true,
  morecomment=[l]{\#},
  morestring=[b]',
  morestring=[b]"
}
\lstdefinelanguage{sql}{
  morekeywords={SELECT,FROM,WHERE,INSERT,INTO,VALUES,UPDATE,SET,DELETE,CREATE,TABLE,PRIMARY,KEY,FOREIGN,REFERENCES,UNIQUE,NOT,NULL,INDEX,ON,ALTER,DROP,CONSTRAINT},
  sensitive=false,
  morecomment=[l]{--}
}
\lstdefinelanguage{json}{
  morecomment=[l]{//},
  morestring=[b]",
  showstringspaces=false
}
\lstdefinelanguage{typescript}{
  morekeywords={import,export,from,const,let,var,function,return,if,else,for,of,in,while,switch,case,break,async,await,try,catch,finally,throw,new,class,this,extends,implements,interface,type,enum,readonly,public,private,protected,static,get,set,true,false,null,undefined,void,string,number,boolean,any,unknown,never,typeof,keyof,as},
  sensitive=true,
  morecomment=[l]{//},
  morecomment=[s]{/*}{*/},
  morestring=[b]',
  morestring=[b]",
  morestring=[b]`
}
\lstalias{javascript}{typescript}
\lstset{
  backgroundcolor=\color{codebg},
  basicstyle=\ttfamily\footnotesize,
  frame=single,
  rulecolor=\color{codeframeclr},
  frameround=tttt,
  breaklines=true,
  tabsize=2,
  showstringspaces=false,
  upquote=true,
  columns=fullflexible
}

% ---------- Helper macros ----------
\newcommand{\codefile}[1]{\texttt{#1}}
\newcommand{\kbd}[1]{\texttt{#1}}

% ---------- Code frame environment ----------
% Automatski dodaje fragile i shrink=5, koristi se umjesto \begin{frame}
% Primjer: \begin{codeframe}{Naslov} ... \end{codeframe}
% Primjer s opcijama: \begin{codeframe}[shrink=15]{Naslov} ... \end{codeframe}
\newenvironment{codeframe}[2][]
  {\begin{frame}[fragile,shrink=5,#1]{#2}}
  {\end{frame}}

% ============================================================

\documentclass[aspectratio=149,11pt]{beamer}
\usetheme{CambridgeUS}

\usepackage{fontspec}              % radi samo sa XeLaTeX/LuaLaTeX
\setsansfont{Latin Modern Sans}    % sans-serif font
\setmonofont{Latin Modern Mono}    % monospace font za kod
\usepackage[croatian]{babel}
\usepackage{csquotes}              % \enquote{} → hrvatski navodnici „..."
\usepackage{lmodern}
\usepackage{booktabs}
\usepackage{listings}
\usepackage{xcolor}
\usepackage{hyperref}

% ---------- Autor i logo ----------
\author{mag.ing. Josip Torić (SIT UNIZD)}
\titlegraphic{\includegraphics[height=2.5cm]{../common/SIT_logo_tekst1.pdf}}

% ---------- Listings ----------
\definecolor{codebg}{RGB}{248,248,248}
\definecolor{codeframeclr}{RGB}{220,220,220}
\lstdefinelanguage{bash}{
  morekeywords={sudo,apt,brew,pip,pipx,uv,winget,python,git,flask,source,deactivate,code,gunicorn,docker,compose,heroku,render,railway,node,npm,pnpm,yarn,uvicorn,alembic,pytest,ruff,black,precommit},
  sensitive=false,
  morecomment=[l]{\#}
}
\lstdefinelanguage{python}{
  morekeywords={from,import,def,class,return,if,elif,else,for,in,while,try,except,raise,with,as,not,and,or,is,lambda,yield,True,False,None,await,async},
  sensitive=true,
  morecomment=[l]{\#},
  morestring=[b]',
  morestring=[b]"
}
\lstdefinelanguage{sql}{
  morekeywords={SELECT,FROM,WHERE,INSERT,INTO,VALUES,UPDATE,SET,DELETE,CREATE,TABLE,PRIMARY,KEY,FOREIGN,REFERENCES,UNIQUE,NOT,NULL,INDEX,ON,ALTER,DROP,CONSTRAINT},
  sensitive=false,
  morecomment=[l]{--}
}
\lstdefinelanguage{json}{
  morecomment=[l]{//},
  morestring=[b]",
  showstringspaces=false
}
\lstdefinelanguage{typescript}{
  morekeywords={import,export,from,const,let,var,function,return,if,else,for,of,in,while,switch,case,break,async,await,try,catch,finally,throw,new,class,this,extends,implements,interface,type,enum,readonly,public,private,protected,static,get,set,true,false,null,undefined,void,string,number,boolean,any,unknown,never,typeof,keyof,as},
  sensitive=true,
  morecomment=[l]{//},
  morecomment=[s]{/*}{*/},
  morestring=[b]',
  morestring=[b]",
  morestring=[b]`
}
\lstalias{javascript}{typescript}
\lstset{
  backgroundcolor=\color{codebg},
  basicstyle=\ttfamily\footnotesize,
  frame=single,
  rulecolor=\color{codeframeclr},
  frameround=tttt,
  breaklines=true,
  tabsize=2,
  showstringspaces=false,
  upquote=true,
  columns=fullflexible
}

% ---------- Helper macros ----------
\newcommand{\codefile}[1]{\texttt{#1}}
\newcommand{\kbd}[1]{\texttt{#1}}

% ---------- Code frame environment ----------
% Automatski dodaje fragile i shrink=5, koristi se umjesto \begin{frame}
% Primjer: \begin{codeframe}{Naslov} ... \end{codeframe}
% Primjer s opcijama: \begin{codeframe}[shrink=15]{Naslov} ... \end{codeframe}
\newenvironment{codeframe}[2][]
  {\begin{frame}[fragile,shrink=5,#1]{#2}}
  {\end{frame}}


\title[Predavanje 0: O predmetu]{Programiranje za web \\ Predavanje 0: O predmetu}
\date{08.10.2026.}

\begin{document}

\frame{\titlepage}

% ------------------------------
\begin{frame}{Plan predavanja}
\begin{enumerate}
  \item Predmet kolegija i ishodi učenja
  \item Referentna aplikacija: prijava klubova na natjecanje
  \item Arhitektura i tehnologije
  \item Organizacija nastave i vrednovanje
  \item Radna okolina i način rada
\end{enumerate}

\vspace{0.6em}
Kolegij je organiziran oko jedne aplikacije koja se razvija od prve do posljednje nastavne jedinice: podatkovni model, poslovna pravila, korisničko sučelje i postavljanje u produkciju.
\end{frame}

% ==============================
\sekcija{Kolegij}{Predmet kolegija, ishodi učenja, pretpostavljeno znanje.}

% ------------------------------
\begin{frame}{Predmet kolegija}
Web aplikacija je program kojemu korisnik pristupa preglednikom, a čiji podatci ostaju pohranjeni i nakon što korisnik zatvori karticu. Za to su nam potrebni:

\begin{itemize}
  \item \textbf{frontend} (klijentski dio): prikaz podataka u pregledniku i unos korisnika
  \item \textbf{backend} (poslužiteljski dio): provjera identiteta i prava, poslovna pravila, pristup podatcima
  \item \textbf{baza podataka}: trajna pohrana i ograničenja nad podatcima
\end{itemize}

\vspace{0.4em}
Kolegij obrađuje sva tri dijela i sučelje kojim međusobno komuniciraju.
\end{frame}

% ------------------------------
\begin{frame}{Ishodi učenja}
Nakon položenog kolegija student može:

\begin{itemize}
  \item modelirati domenu i preslikati je u relacijsku bazu zajedno s ograničenjima
  \item oblikovati programsko sučelje API (engl. \emph{application programming interface}) u stilu REST (engl. \emph{representational state transfer}) i validirati ulazne podatke
  \item implementirati autentikaciju tokenom JWT (engl. \emph{JSON Web Token}) te autorizaciju po ulozi i po vlasništvu zapisa
  \item izraditi jednostraničnu aplikaciju (engl. \emph{single-page application}, SPA) koja se oslanja na taj API
  \item napisati testove poslovnih pravila i postaviti aplikaciju u produkcijsku okolinu
\end{itemize}
\end{frame}

% ------------------------------
\begin{frame}{Pretpostavljeno znanje}
\begin{itemize}
  \item programiranje u jeziku opće namjene; sintaksa Pythona uvodi se usput
  \item osnove relacijskih baza: tablica, primarni i strani ključ, upit \texttt{SELECT}
  \item osnove HTML-a i CSS-a
  \item rad s Gitom: \texttt{commit}, \texttt{branch}, \texttt{remote}
\end{itemize}

\vspace{0.4em}

\end{frame}

% ==============================
\sekcija{Referentna aplikacija}{Zadatak, model domene, faze i rokovi, svojstva domene relevantna za kolegij.}

% ------------------------------
\begin{frame}{Zadatak: aplikacija za prijavu na natjecanje}
Organizator powerlifting natjecanja prikuplja prijave od klubova. Bez aplikacije prijave stižu e-poštom i tablicama u Excelu, pa organizator ručno provjerava:

\begin{itemize}
  \item je li prijava stigla prije isteka roka
  \item prijavljuje li klub natjecatelja koji mu pripada
  \item je li isti natjecatelj prijavljen dvaput
  \item odgovara li težinska kategorija spolu i tjelesnoj masi natjecatelja
\end{itemize}

\vspace{0.4em}
Svaka od tih provjera je pravilo koje aplikacija treba provoditi umjesto čovjeka. Ista se aplikacija razvija kroz cijeli kolegij i služi kao predložak za studentski projekt.
\end{frame}

% ------------------------------
\begin{frame}{Model domene}
\begin{columns}[T]
\column{0.5\textwidth}
\textbf{Entiteti}
\begin{itemize}
  \item \textbf{Klub}: prijavljuje natjecatelje
  \item \textbf{Korisnik}: račun za prijavu, uloga \texttt{admin} ili \texttt{club}; svaki klub ima jedan
  \item \textbf{Natjecatelj}: pripada točno jednom klubu
  \item \textbf{Natjecanje}: datum održavanja i dva roka
  \item \textbf{Prijava}: veza natjecatelja i natjecanja
\end{itemize}

\column{0.5\textwidth}
\textbf{Pravila nad modelom}
\begin{itemize}
  \item klub vidi i mijenja samo svoje natjecatelje i svoje prijave
  \item organizator vidi prijave svih klubova
  \item isti natjecatelj ne može imati dvije aktivne prijave na istom natjecanju
  \item težinska kategorija bira se iz popisa kategorija saveza IPF
\end{itemize}
\end{columns}

\vspace{0.5em}
{\small Kategorije IPF, muškarci: 59, 66, 74, 83, 93, 105, 120 i 120+ kg; žene: 47, 52, 57, 63, 69, 76, 84 i 84+ kg.}
\end{frame}

% ------------------------------
\begin{frame}{Faze natjecanja i rokovi}
Natjecanje ima dva roka nominacije: \textbf{preliminarni} (\codefile{prelim\_deadline}) i \textbf{finalni} (\codefile{final\_deadline}). Oni dijele vrijeme na tri faze s različitim skupom dopuštenih operacija.

\vspace{0.5em}
{\small
\begin{tabular}{@{}llccc@{}}
  \toprule
  \textbf{Faza} & \textbf{Vrijeme} & \textbf{Nova prijava} & \textbf{Promjena kat.} & \textbf{Odjava} \\
  \midrule
  \texttt{OPEN}           & do preliminarnog roka & da & da & da \\
  \texttt{PRELIM\_PASSED} & između dvaju rokova   & ne & da & da \\
  \texttt{CLOSED}         & nakon finalnog roka   & ne & ne & ne \\
  \bottomrule
\end{tabular}
}

\vspace{0.5em}
Preliminarni rok organizatoru daje popis potreban za raspored vaganja i natjecateljskih skupina; nakon finalnog roka popis se više ne mijenja.
\end{frame}

% ------------------------------
\begin{frame}{Radni primjer: četiri zahtjeva istoga kluba}
Natjecanje se održava 14.03.2027.; preliminarni rok istječe 14.02.2027. u 23:59 (28 dana ranije), finalni 07.03.2027. u 23:59 (7 dana ranije).

\vspace{0.5em}
{\footnotesize
\setlength{\tabcolsep}{4pt}
\begin{tabular}{@{}llll@{}}
  \toprule
  \textbf{Trenutak zahtjeva} & \textbf{Zahtjev} & \textbf{Faza} & \textbf{Odgovor} \\
  \midrule
  10.02.2027. & prijava natjecatelja, kat. 83 kg   & \texttt{OPEN}           & \texttt{201 Created} \\
  20.02.2027. & promjena kategorije 83 $\to$ 93 kg & \texttt{PRELIM\_PASSED} & \texttt{200 OK} \\
  20.02.2027. & prijava novog natjecatelja         & \texttt{PRELIM\_PASSED} & \texttt{400 Bad Request} \\
  09.03.2027. & odjava natjecatelja                & \texttt{CLOSED}         & \texttt{400 Bad Request} \\
  \bottomrule
\end{tabular}
}

\vspace{0.5em}
\end{frame}

% ------------------------------
\begin{frame}{Svojstva domene relevantna za kolegij}
\begin{itemize}
  \item postoje dvije uloge, pa se autorizacija provjerava pri svakom zahtjevu
  \item vlasništvo nad zapisom: klub A ne smije dohvatiti natjecatelja kluba B
  \item pravila ovisna o vremenu: isti zahtjev prolazi ili ne prolazi ovisno o fazi
  \item relacije i ograničenja jedinstvenosti, koja provodi baza, a ne aplikacijski kod
  \item sučelje s tablicama, filtrima i dostupnošću naredbi ovisnom o fazi
\end{itemize}

\vspace{0.4em}
\end{frame}

% ==============================
\sekcija{Arhitektura i tehnologije}{Troslojna arhitektura, podjela odgovornosti, tehnologije, struktura repozitorija.}

% ------------------------------
\begin{frame}{Troslojna arhitektura}
\begin{center}
\begin{tikzpicture}[
  node distance=6mm and 20mm,
  box/.style={draw=ink, thick, rounded corners, fill=white, align=center,
              minimum height=12mm, minimum width=24mm, font=\small, text=textclr},
  lbl/.style={font=\scriptsize, align=center, text=mutedclr}
]
  \node[box] (fe) {Preglednik \\ (SPA)};
  \node[box, right=of fe] (be) {Poslužitelj \\ (API)};
  \node[box, right=of be] (db) {Baza \\ (PostgreSQL)};

  \draw[-{Latex}, ink2, thick] ([yshift=3mm]fe.east) -- node[lbl, above] {zahtjev \\ JSON, token} ([yshift=3mm]be.west);
  \draw[-{Latex}, ink2, thick] ([yshift=-3mm]be.west) -- node[lbl, below] {odgovor \\ JSON, status} ([yshift=-3mm]fe.east);
  \draw[-{Latex}, ink2, thick] ([yshift=3mm]be.east) -- node[lbl, above] {upit \\ SQL} ([yshift=3mm]db.west);
  \draw[-{Latex}, ink2, thick] ([yshift=-3mm]db.west) -- node[lbl, below] {zapisi} ([yshift=-3mm]be.east);
\end{tikzpicture}
\end{center}

\vspace{0.3em}
Preglednik ne pristupa bazi izravno. Svaki podatak koji korisnik vidi prošao je kroz \emph{backend}, koji je pritom provjerio tko je korisnik i što smije.
\end{frame}

% ------------------------------
\begin{frame}{Podjela odgovornosti}
\begin{itemize}
  \item Baza je izvor istine o podatcima; ograničenja zapisana u shemi provode se neovisno o tome koji ih kod pokušava prekršiti.
  \item Backend je mjesto na kojem se pravilo provodi: prije svake izmjene provjeravaju se identitet, uloga, vlasništvo nad zapisom i faza natjecanja.
  \item \emph{Frontend} pravila prikazuje, ali ih ne provodi. Onemogućen gumb korisniku štedi uzaludan pokušaj, ali ne sprječava zahtjev poslan mimo sučelja.
\end{itemize}

\vspace{0.4em}
Provjera ispravnosti te podjele: zahtjev poslan izravno na API, mimo korisničkog sučelja, mora dobiti isti odgovor kao zahtjev poslan iz preglednika. Ako ga ne dobije, pravilo je zapisano na krivome mjestu.
\end{frame}

% ------------------------------
\begin{frame}{Tehnologije}
\begin{block}{Backend}
FastAPI i Uvicorn (HTTP) \;|\; Pydantic (validacija) \;|\; SQLAlchemy 2.0 (ORM) \;|\; Alembic (migracije) \;|\; PostgreSQL
\end{block}

\begin{block}{Frontend}
Vue 3, Vite i TypeScript \;|\; Vue Router (rute) \;|\; Pinia (stanje) \;|\; Axios (HTTP)
\end{block}

\begin{block}{Alati i okolina}
Docker Compose (baza) \;|\; Git i GitHub Actions (CI) \;|\; uv (Python okolina) \;|\; pytest (testovi) \;|\; ruff (statička analiza) \;|\; Railway ili Render (postavljanje u rad)
\end{block}

\vspace{0.2em}
Ovo je jedna moguća kombinacija tehnologija. Ista se načela primjenjuju i uz drukčiji izbor.
\end{frame}

% ------------------------------
\begin{frame}{Struktura repozitorija}
\begin{columns}[T]
\column{0.48\textwidth}
\textbf{Backend}
\begin{itemize}
  \item \codefile{routers/}: rute i parametri zahtjeva
  \item \codefile{schemas/}: Pydantic modeli ulaza i izlaza
  \item \codefile{models/}: SQLAlchemy tablice
  \item \codefile{services/}: poslovna pravila
  \item \codefile{repositories/}: upiti nad bazom
  \item \codefile{core/}: konfiguracija, autentikacija, pogreške
\end{itemize}

\column{0.48\textwidth}
\textbf{Frontend}
\begin{itemize}
  \item \codefile{src/views/}: stranice
  \item \codefile{src/components/}: komponente sučelja
  \item \codefile{src/stores/}: spremišta stanja (Pinia)
  \item \codefile{src/router/}: rute i njihove zaštite
  \item \codefile{src/api/}: omotač oko HTTP klijenta
\end{itemize}
\end{columns}

\vspace{0.2em}
\end{frame}

% ==============================
\sekcija{Organizacija nastave}{Raspored predavanja, zahtjevi na studentski projekt, bodovanje, obrana, akademska čestitost.}

% ------------------------------
\begin{frame}{Raspored predavanja}
{\footnotesize
	\renewcommand{\arraystretch}{1.0}
	\begin{tabular}{@{}lll@{}}
		\toprule
		\textbf{Pred.} & \textbf{Datum} & \textbf{Tema} \\
		\midrule
		0  & 08.10. & O predmetu \\
		1  & 08.10. & Arhitektura backenda i kostur projekta \\
		2  & 15.10. & Modeliranje podataka, ORM i migracije \\
		3  & 22.10. & Autentikacija i sigurnost (JWT) \\
		4  & 29.10. & Autorizacija i vlasništvo nad zapisom \\
		5  & 05.11. & Oblikovanje API-ja i CRUD domene \\
		6  & 12.11. & Poslovna pravila i faze u backendu \\
		   & 19.11. & \textbf{Obrana prvog dijela projekta} \\
		7  & 26.11. & Arhitektura jednostranične aplikacije (Vue) \\
		8  & 03.12. & Komunikacija s API-jem i tijek prijave korisnika \\
		9  & 10.12. & Upravljanje stanjem i stanja korisničkog sučelja \\
		10 & 17.12. & Tablice, filtri i administratorsko sučelje \\
		   & 24.12. i 31.12. & \emph{Božićni i novogodišnji praznici} \\
		11 & 07.01. & Testiranje cjelovite aplikacije \\
		12 & 14.01. & Produkcijska izgradnja i postavljanje u produkciju \\
		   & 21.01. & \textbf{Obrana drugog dijela projekta} \\
		\bottomrule
	\end{tabular}
}
\end{frame}

% ------------------------------
\begin{frame}{Studentski projekt}
Tijekom semestra svaki student izrađuje vlastitu aplikaciju, po uzoru na referentnu, ali u drugoj domeni. Zahtjevi:

\begin{itemize}
  \item autentikacija tokenom i najmanje dvije uloge s različitim pravima
  \item najmanje dva entiteta povezana relacijom
  \item vlasništvo nad zapisom: korisnik vidi i mijenja samo svoje podatke
  \item najmanje jedno pravilo ovisno o stanju ili vremenu: status, faza ili rok
  \item sučelje sa stanjima učitavanja, pogreške i praznog skupa rezultata
\end{itemize}

\vspace{0.4em}
Predaja je u dva dijela, svaki nosi 40 bodova: prvi obuhvaća backend, drugi cjelovitu aplikaciju postavljenu u rad.
\end{frame}

% ------------------------------
\begin{frame}{Primjeri tema projekta}
\begin{itemize}
  \item rezervacija termina: dvorane, treneri, konzultacije
  \item prijave na događanja: kapacitet, lista čekanja, evidencija dolaska
  \item praćenje zadataka: projekti, statusi, komentari
  \item planovi treninga: programi, izvedbe, napredak
  \item odobravanje dokumenata: zahtjev, razine odobrenja, verzije
  \item mali oglasnik: oglasi, poruke, spremljeni oglasi
  \item evidencija članstva: članovi, uplate, statusi
\end{itemize}

\vspace{0.5em}
Tema mora sadržavati uloge i pravila. Aplikacija koja zapise samo unosi i prikazuje, bez pravila o tome tko što smije i kada, ne zadovoljava zahtjeve projekta.
\end{frame}

% ------------------------------
\begin{frame}{Bodovanje}
\begin{tabular}{@{}lrr@{}}
  \toprule
  \textbf{Komponenta} & \textbf{Bodovi} & \textbf{Prag} \\
  \midrule
  Prvi dio projekta (backend) & 40 & 15 \\
  Drugi dio projekta (cjelovita aplikacija u radu) & 40 & 15 \\
  Usmeni ispit & 20 & 10 \\
  \midrule
  \textbf{Ukupno} & \textbf{100} & \\
  \bottomrule
\end{tabular}

\vspace{0.6em}
\begin{center}
\begin{tabular}{@{}lr@{}}
  \toprule
  \textbf{Bodovi} & \textbf{Ocjena} \\
  \midrule
  50 do 59   & dovoljan (2) \\
  60 do 74   & dobar (3) \\
  75 do 87   & vrlo dobar (4) \\
  88 do 100  & izvrstan (5) \\
  \bottomrule
\end{tabular}
\end{center}
\end{frame}

% ------------------------------
\begin{frame}{Obrana projekta}
Obrana se sastoji od prikaza rada aplikacije, pregleda repozitorija i objašnjenja dijelova koda koje odabere nastavnik. Moguća je i manja izmjena uživo: dodavanje polja, migracija te prilagodba modela i sučelja.

\begin{block}{Primjer pitanja na obrani projekta}
\begin{itemize}
  \item Na kojem se mjestu provjerava vlasništvo nad zapisom?
  \item Kako backend odbija izmjenu nakon isteka roka?
  \item Zašto zahtjev vraća \texttt{401}, a ne \texttt{403}? Zašto \texttt{409}, a ne \texttt{422}?
  \item Gdje se na frontendu obrađuje odgovor \texttt{401}?
\end{itemize}
\end{block}

Ocjenjuje se objašnjenje odluka ugrađenih u kod.
\end{frame}

% ------------------------------
\begin{frame}{Akademska čestitost i uporaba AI alata}
Uporaba AI alata i vanjskih izvora dopuštena je. Uvjet je da student razumije i objasni svaki dio predanoga koda; na obrani se provjerava razumijevanje.

\vspace{0.4em}
Za svaki preuzeti isječak koda treba znati:
\begin{itemize}
  \item koji problem rješava i uz koje pretpostavke
  \item kako se ponaša u rubnim slučajevima: prazan ulaz, istekao \emph{token}, istodobni zahtjevi
  \item koji ga test pokriva
\end{itemize}

\vspace{0.3em}
\end{frame}

% ==============================
\sekcija{Radna okolina}{Potrebni alati, pokretanje projekta, kriteriji kvalitete koda, način rada.}

% ------------------------------
\begin{frame}{Potrebni alati}
\begin{itemize}
  \item Git
  \item Python 3.11 ili noviji
  \item uv, za virtualnu okolinu i instalaciju paketa
  \item Docker Desktop s alatom Docker Compose, za pokretanje baze
  \item Node.js 18 ili noviji
  \item uređivač koda po izboru; na predavanjima se koristi VS Code
\end{itemize}

\vspace{0.5em}
Okolina se postavlja prije prvoga predavanja s kodom.
\end{frame}

% ------------------------------
\begin{frame}[fragile]{Pokretanje projekta: baza i backend}
\begin{lstlisting}[language=bash]
# 1) Environment variables
cp .env.example .env         # PowerShell: Copy-Item .env.example .env

# 2) Database (PostgreSQL in a container)
docker compose up -d db

# 3) Backend
cd api
uv venv
source .venv/bin/activate    # PowerShell: .venv\Scripts\Activate.ps1
uv pip install -r requirements.txt
alembic upgrade head         # apply migrations
python -m app.seed           # seed data
uvicorn app.main:app --reload
\end{lstlisting}

\vspace{0.3em}
\codefile{.env.example} je u repozitoriju i služi kao popis varijabli; \codefile{.env} s vlastitim vrijednostima ostaje na računalu i ne ulazi u repozitorij.
\end{frame}

% ------------------------------
\begin{frame}[fragile]{Pokretanje projekta: frontend i provjera}
\begin{lstlisting}[language=bash]
# 4) Frontend
cd ../web
npm install
npm run dev
\end{lstlisting}

\vspace{0.5em}
Nakon pokretanja backend odgovara na \texttt{/health}, a opis sučelja stoji na \texttt{/docs}. Sučelje je dostupno na adresi koju ispiše \texttt{npm run dev}.

\end{frame}

% ------------------------------
\begin{frame}{Kriteriji kvalitete koda}
\begin{columns}[T]
\column{0.48\textwidth}
\textbf{Backend}
\begin{itemize}
  \item validacija svih ulaznih podataka
  \item jedinstven oblik odgovora s pogreškom
  \item pravila u sloju \codefile{services/}, a ne u ruti
  \item test za svako poslovno pravilo
\end{itemize}

\column{0.48\textwidth}
\textbf{Frontend}
\begin{itemize}
  \item stanja učitavanja, pogreške i praznog skupa
  \item zaštita ruta prema ulozi korisnika
  \item validacija obrazaca prije slanja
  \item adresa API-ja iz varijable okoline
\end{itemize}
\end{columns}

\vspace{0.5em}
Objema stranama zajedničko je načelo: svako se pravilo zapisuje na jednome mjestu i odatle poziva, umjesto da se ponavlja u svakoj ruti ili komponenti.
\end{frame}

% ------------------------------
\begin{frame}{Način rada}
\begin{itemize}
  \item Nakon svakoga predavanja: ponoviti korake u vlastitom repozitoriju i usporediti razlike naredbom \kbd{git diff}.
  \item Kratak i čest rad daje bolji rezultat od dugoga rada neposredno prije roka predaje.
  \item Postupak pri pogrešci: pročitati poruku, izolirati najmanji primjer koji je izaziva, provjeriti zahtjev i odgovor (statusni kod i tijelo), pa tek onda mijenjati kod.
  \item Pitanje nastavnom osoblju postavlja se uz konkretan materijal: poruku pogreške te sadržaj zahtjeva i odgovora.
\end{itemize}
\end{frame}

% ------------------------------
\begin{frame}{Sažetak}
\begin{itemize}
  \item Web aplikacija ima tri dijela; pravila provodi backend, a baza je izvor istine o podatcima.
  \item Referentna domena su prijave klubova na natjecanje: dvije uloge, vlasništvo nad zapisom te tri faze određene dvama rokovima.
  \item Faza određuje ishod zahtjeva: nova prijava prolazi samo u fazi \texttt{OPEN}, promjena kategorije i odjava i u fazi \texttt{PRELIM\_PASSED}, a nakon finalnoga roka nijedna izmjena.
  \item Studentski projekt nosi 80 od 100 bodova i predaje se u dva dijela, a usmeni ispit preostalih 20.
  \item Na obrani se ocjenjuje objašnjenje odluka ugrađenih u kod.
\end{itemize}
\end{frame}

% ------------------------------
\begin{frame}{Što slijedi}
\begin{itemize}
  \item \textbf{08.10.2026.} Predavanje 1, arhitektura backenda i kostur projekta: struktura repozitorija, baza u kontejneru, prva ruta API-ja.
  \item \textbf{15.10.2026.} Predavanje 2, modeliranje podataka, ORM i migracije.
  \item \textbf{19.11.2026.} Obrana prvoga dijela projekta.
\end{itemize}

\vspace{0.5em}
\end{frame}

\end{document}
